\documentclass[11pt,a4paper]{article}

\usepackage[T1]{fontenc}
\usepackage[utf8]{inputenc}
\usepackage[a4paper,margin=25mm]{geometry}
\usepackage{xcolor}
\usepackage{array}
\usepackage{booktabs}
\usepackage{tabularx}
\usepackage{listings}
\usepackage{graphicx}
\usepackage{float}
\usepackage{hyperref}

\hypersetup{
    colorlinks=true,
    linkcolor=blue,
    urlcolor=blue
}

\definecolor{codegray}{gray}{0.95}

\lstset{
    backgroundcolor=\color{codegray},
    basicstyle=\ttfamily\small,
    frame=single,
    breaklines=true,
    columns=fullflexible,
    keepspaces=true,
    showstringspaces=false,
    tabsize=4
}

\newcommand{\code}[1]{\texttt{#1}}
% \shot{file under results/screenshots/}{width}{caption}
\newcommand{\shot}[3]{%
    \begin{figure}[H]
        \centering
        \includegraphics[width=#2,height=0.6\textheight,keepaspectratio]{results/screenshots/#1}
        \caption{#3}
    \end{figure}
}

\title{CPE333 Problem Session 12\\File System and Implementation}
\author{
    Guntee Doungmanee (ID: 67070501003)\\
    Drarun Yookphan (ID: 67070501016)\\
    Rattathamnoon Buranapattana (ID: 67070501037)\\
    Supawit Marayat (ID: 67070501045)
}
\date{}

\begin{document}

\maketitle
\tableofcontents
\newpage

\section{Introduction}

This report covers the three experiments of Problem Session 12. Lab 1 builds a
general graph directory structure with hard links, symbolic links, and
directory links on Linux and Windows. Lab 2 studies the Linux named pipe (FIFO)
file type. Lab 3 traces the system calls that \code{cat} uses to read a file and
write it to the screen.

The Linux transcripts were produced by \code{run-linux.sh}, and the Windows
transcript by \code{windows-links.cmd} (which calls \code{run-windows.ps1}).
Every listing in this report is copied from \code{results/}, and every figure
is from \code{results/screenshots/}. Inode numbers, dates, and free space change
from run to run.

\section{Execution environment}

\paragraph{Linux.} Labs 1--3 ran directly on CachyOS, an Arch-based
distribution, on an x86-64 machine.

\begin{lstlisting}
$ uname -srm
Linux 7.3.3-1-cachyos x86_64
$ grep PRETTY_NAME /etc/os-release
PRETTY_NAME="CachyOS"
$ ls --version | head -1
ls (GNU coreutils) 9.12
$ strace --version | head -1
strace -- version 7.3
$ tmux -V
tmux 3.8
$ date
Sun Nov 22 06:01:05 PM +07 2026
\end{lstlisting}

\shot{linux/00_environment.png}{0.6\linewidth}{Linux environment.}

\paragraph{Windows.} Lab 1 ran on native Windows 11 from an Administrator
Command Prompt. The machine has two separate NTFS volumes, \code{C:} and
\code{D:}, which are needed for the cross-drive test.

\begin{lstlisting}
> ver
Microsoft Windows [Version 10.0.22631.4751]
> fsutil fsinfo volumeinfo C:
Volume Name : Windows-SSD
File System Name : NTFS
Supports Reparse Points
Supports Hard Links
...
> fsutil fsinfo volumeinfo D:
Volume Name : New Volume
File System Name : NTFS
Supports Reparse Points
Supports Hard Links
...
\end{lstlisting}

The Windows username, the temporary run identifier, and the volume serial
numbers are replaced with \code{USER}, \code{RUNID}, and \code{[REDACTED]} in
the transcript and covered with gray boxes in the screenshots. The Windows
screenshots show a read-only viewer of the captured command output, with two
columns read from left to right.

\shot{windows/00_environment.png}{\linewidth}{Windows environment and NTFS volumes.}

\section{Lab 1: Symbolic and hard links}

\subsection{Background: general graph directory}

In a tree-structured directory, each file has exactly one path. A link adds a
second path to an existing file or directory, so the directory structure
becomes a graph. If links may point to directories, the graph can contain a
cycle. The operating system must then handle three problems: a file can be
reached by several names, deletion must not leave dangling references, and a
recursive traversal must not loop forever.

Both systems offer two kinds of link:

\begin{description}
    \item[Hard link] A second directory entry that points to the same inode
    (Linux) or the same MFT file record (NTFS). All names are equal; none is the
    ``original''. The file system keeps a link count and frees the data only
    when the count reaches zero.
    \item[Symbolic (soft) link] A small separate file that stores a path. The
    kernel follows the path each time the link is opened. The target may be
    missing, on another file system, or a directory.
\end{description}

Windows adds a third kind, the \emph{junction} (\code{mklink /J}), which is a
directory-only reparse point that stores an absolute local path.

\begin{center}
\begin{tabular}{@{}lll@{}}
\toprule
Link type & Linux & Windows (\code{cmd.exe}) \\
\midrule
Hard link to file & \code{ln src link} & \code{mklink /H link src} \\
Symbolic link to file & \code{ln -s src link} & \code{mklink link src} \\
Symbolic link to directory & \code{ln -s dir link} & \code{mklink /D link dir} \\
Junction to directory & -- & \code{mklink /J link dir} \\
\bottomrule
\end{tabular}
\end{center}

\code{ln} takes the source first, while \code{mklink} takes the new link name
first.

\subsection{Hard and soft links, then delete the source}

\paragraph{Linux.} Create a source file, a hard link, and a soft link, then list
them with inode numbers:

\begin{lstlisting}
$ echo "original data" > source.txt
$ ln source.txt hard.txt
$ ln -s source.txt soft.txt
$ ls -li
total 8
3016 -rw-r--r-- 2 gunteed gunteed 14 Nov 22 18:01 hard.txt
3017 lrwxrwxrwx 1 gunteed gunteed 10 Nov 22 18:01 soft.txt -> source.txt
3016 -rw-r--r-- 2 gunteed gunteed 14 Nov 22 18:01 source.txt
\end{lstlisting}

\code{hard.txt} and \code{source.txt} share inode 3016, and the link count
column shows 2. They are two names for one file. \code{soft.txt} has its own
inode (3017), type \code{l}, and a size of 10 bytes, which is the length of the
string \code{source.txt} that it stores.

Delete the source:

\begin{lstlisting}
$ rm source.txt
$ ls -li
total 4
3016 -rw-r--r-- 1 gunteed gunteed 14 Nov 22 18:01 hard.txt
3017 lrwxrwxrwx 1 gunteed gunteed 10 Nov 22 18:01 soft.txt -> source.txt
$ cat hard.txt
original data
$ cat soft.txt
cat: soft.txt: No such file or directory
$ readlink soft.txt
source.txt
\end{lstlisting}

\code{rm} removes one directory entry and lowers the link count from 2 to 1.
The inode and its data stay because \code{hard.txt} still refers to them, so
\code{cat hard.txt} still works. The symbolic link still exists and still
stores the path \code{source.txt}, but that path no longer resolves. It is now
a \emph{dangling} link.

Create a new file with the same name:

\begin{lstlisting}
$ echo "new source" > source.txt
$ cat soft.txt
new source
$ cat hard.txt
original data
\end{lstlisting}

The soft link works again, but it now shows the new file, because it follows a
name and not an inode. The hard link keeps the old data, because it is tied to
the old inode.

\shot{linux/01_lab1_basic_links.png}{0.7\linewidth}{Linux: hard and soft link
before and after deleting the source.}

\paragraph{Windows.}\mbox{}

\begin{lstlisting}
> echo original data> source.txt
> mklink /H hard.txt source.txt
Hardlink created for hard.txt <<===>> source.txt
> mklink soft.txt source.txt
symbolic link created for soft.txt <<===>> source.txt
> fsutil hardlink list hard.txt
\Users\USER\AppData\Local\Temp\ps12-RUNID\source.txt
\Users\USER\AppData\Local\Temp\ps12-RUNID\hard.txt
> del source.txt
> type hard.txt
original data
> type soft.txt
The system cannot find the file specified.
> echo new source> source.txt
> type soft.txt
new source
> type hard.txt
original data
\end{lstlisting}

NTFS behaves the same way as Linux. \code{fsutil hardlink list} shows two names
for one file record. After \code{del source.txt}, the hard link still holds the
data, and the symbolic link fails until a new \code{source.txt} is created.

\shot{windows/01_lab1_basic_links.png}{\linewidth}{Windows: hard and symbolic
link before and after deleting the source.}

\paragraph{Answer.} After the source is deleted, the hard link still works and
holds the full data, because deleting a name only lowers the link count. The
symbolic link remains as a file but is broken, because the path it stores no
longer exists. This is the same on Linux and Windows.

\subsection{Level 1 and level 2 links}

A level 1 link points to the source, and a level 2 link points to level 1.

\paragraph{Symbolic chain, delete the source (Linux).}\mbox{}

\begin{lstlisting}
$ ln -s source.txt level1
$ ln -s level1 level2
$ cat level2
level 0 data
$ rm source.txt
$ cat level1
cat: level1: No such file or directory
$ cat level2
cat: level2: No such file or directory
\end{lstlisting}

The kernel resolves \code{level2} to \code{level1}, then to
\code{source.txt}. When the end of the chain is missing, every link in the
chain is broken.

\paragraph{Symbolic chain, delete level 1 (Linux).}\mbox{}

\begin{lstlisting}
$ rm level1
$ ls -li
total 4
3026 lrwxrwxrwx 1 gunteed gunteed  6 Nov 22 18:01 level2 -> level1
3024 -rw-r--r-- 1 gunteed gunteed 13 Nov 22 18:01 source.txt
$ cat source.txt
level 0 data
$ cat level2
cat: level2: No such file or directory
\end{lstlisting}

The source is safe, but \code{level2} is broken even though the source still
exists. \code{level2} only knows the name \code{level1}; it does not know the
source.

\shot{linux/02a_lab1_symbolic_link_chain.png}{0.75\linewidth}{Linux: symbolic
link chains.}

\paragraph{Hard chain, delete the source (Linux).}\mbox{}

\begin{lstlisting}
$ ln source.txt level1
$ ln level1 level2
$ ls -li
total 12
3028 -rw-r--r-- 3 gunteed gunteed 13 Nov 22 18:01 level1
3028 -rw-r--r-- 3 gunteed gunteed 13 Nov 22 18:01 level2
3028 -rw-r--r-- 3 gunteed gunteed 13 Nov 22 18:01 source.txt
$ rm source.txt
$ ls -li
total 8
3028 -rw-r--r-- 2 gunteed gunteed 13 Nov 22 18:01 level1
3028 -rw-r--r-- 2 gunteed gunteed 13 Nov 22 18:01 level2
$ cat level2
level 0 data
\end{lstlisting}

\paragraph{Hard chain, delete level 1 (Linux).}\mbox{}

\begin{lstlisting}
$ rm level1
$ ls -li
total 8
3030 -rw-r--r-- 2 gunteed gunteed 13 Nov 22 18:01 level2
3030 -rw-r--r-- 2 gunteed gunteed 13 Nov 22 18:01 source.txt
$ cat level2
level 0 data
\end{lstlisting}

All three hard names share one inode. \code{ln level1 level2} does not create
a link ``to level 1''; it resolves \code{level1} to its inode and adds a third
name for that inode. A hard link chain has no levels. Deleting any one name
only lowers the count from 3 to 2, and the other names keep working.

\shot{linux/02b_lab1_hard_link_chains.png}{0.7\linewidth}{Linux: hard link
chains.}

\paragraph{Windows.} The same four cases give the same results:

\begin{lstlisting}
sym-a> del source.txt
sym-a> type level1
The system cannot find the file specified.
sym-a> type level2
The system cannot find the file specified.

sym-b> del level1
sym-b> type source.txt
level 0 data
sym-b> type level2
The system cannot find the file specified.

hard-a> fsutil hardlink list level2
...\hard-a\level1
...\hard-a\source.txt
...\hard-a\level2
hard-a> del source.txt
hard-a> type level2
level 0 data

hard-b> del level1
hard-b> fsutil hardlink list level2
...\hard-b\level2
...\hard-b\source.txt
hard-b> type level2
level 0 data
\end{lstlisting}

\code{dir /a:l} lists the symbolic links as \code{<SYMLINK> level2 [level1]},
so Windows also stores only the name of the next link. \code{fsutil hardlink
list} shows three names for one file, which drop to two after any one is
deleted.

\shot{windows/02a_lab1_symbolic_delete_source.png}{\linewidth}{Windows: symbolic
chain, delete the source.}
\shot{windows/02b_lab1_symbolic_delete_level1.png}{\linewidth}{Windows: symbolic
chain, delete level 1.}
\shot{windows/02c_lab1_hard_link_chains.png}{\linewidth}{Windows: hard link
chains, delete the source or level 1.}

\paragraph{Answer.} With symbolic links, deleting the source breaks both
levels, and deleting level 1 breaks level 2 even though the source still
exists. With hard links, deleting the source or level 1 does not affect
level 2, because all names point directly to the same file.

\subsection{Links to directories and a loop}

\paragraph{Linux.}\mbox{}

\begin{lstlisting}
$ mkdir -p top/sub
$ echo "file in sub" > top/sub/file.txt
$ ln top top-hard
ln: top: hard link not allowed for directory
$ ln -s top top-soft
$ ls -l top-soft/sub
total 4
-rw-r--r-- 1 gunteed gunteed 12 Nov 22 18:01 file.txt
\end{lstlisting}

Linux refuses hard links to directories (\code{EPERM}). A directory hard link
could create a cycle that the link count cannot detect, so a deleted loop of
directories might never reach a count of zero. Only \code{.} and \code{..} are
hard links to directories, and the file system creates those itself. A
symbolic link to a directory is allowed.

Now create a symbolic link inside \code{top/sub} that points back up to
\code{top}:

\begin{lstlisting}
$ ln -s .. top/sub/up
$ ls -l top/sub
total 4
-rw-r--r-- 1 gunteed gunteed 12 Nov 22 18:01 file.txt
lrwxrwxrwx 1 gunteed gunteed  2 Nov 22 18:01 up -> ..
$ ls top/sub/up/sub/up/sub/up/sub
file.txt
up
$ realpath top/sub/up/sub/up/sub/up
/tmp/ps12.1arIAs/lab1/top
\end{lstlisting}

The graph now has a cycle \code{top} $\rightarrow$ \code{sub} $\rightarrow$
\code{up} $\rightarrow$ \code{top}, so a path can be made arbitrarily long.
\code{realpath} shows that the long path is just \code{top}.

List it recursively:

\begin{lstlisting}
$ ls -R top
top:
sub

top/sub:
file.txt
up
$ ls -RL top
top:
sub

top/sub:
file.txt
up
ls: top/sub/up: not listing already-listed directory
$ find -L top
top
top/sub
find: File system loop detected; the following directory is part of the cycle: 'top/sub/up'
top/sub/file.txt
\end{lstlisting}

\code{ls -R} does not loop. By default it does not follow symbolic links while
recursing, so it prints \code{up} as an entry and stops. When it is told to
follow links (\code{-L}), \code{ls} remembers the device and inode of each
directory it has listed and refuses to list one twice. \code{find -L} uses the
same check and reports a file system loop. The protection is in the tools, not
in the kernel; a naive program that follows symbolic links would recurse until
it hits the kernel's limits on path length or link depth.

\shot{linux/03_lab1_directory_loop.png}{0.85\linewidth}{Linux: directory links
and the loop back to \code{top}.}

\paragraph{Windows.}\mbox{}

\begin{lstlisting}
> mkdir top\sub
> echo file in sub> top\sub\file.txt
> mklink /H top-hard top
Access is denied.
> mklink /D top-soft top
symbolic link created for top-soft <<===>> top
> mklink /J top-junction "C:\Users\USER\AppData\Local\Temp\ps12-RUNID\top"
Junction created for top-junction <<===>> C:\Users\USER\AppData\Local\Temp\ps12-RUNID\top
> type top-soft\sub\file.txt
file in sub
> type top-junction\sub\file.txt
file in sub
\end{lstlisting}

NTFS also forbids directory hard links (\code{Access is denied.}). A directory
symbolic link (\code{<SYMLINKD>}) and a junction (\code{<JUNCTION>}) both work.

\shot{windows/03a_lab1_directory_links.png}{\linewidth}{Windows: directory hard
link, symbolic link, and junction.}

Now create a junction inside \code{top\textbackslash sub} that points back to
\code{top}, and list it recursively:

\begin{lstlisting}
> mklink /J top\sub\up "C:\Users\USER\AppData\Local\Temp\ps12-RUNID\top"
Junction created for top\sub\up <<===>> C:\Users\USER\AppData\Local\Temp\ps12-RUNID\top
> type top\sub\up\sub\up\sub\file.txt
file in sub
> dir /s /b top > loop-list.txt
> find /c /v "" loop-list.txt
---------- LOOP-LIST.TXT: 192
> powershell ... (first 3 paths, last 3 paths, longest path)
First 3 paths:
C:\Users\USER\AppData\Local\Temp\ps12-RUNID\top\sub
C:\Users\USER\AppData\Local\Temp\ps12-RUNID\top\sub\file.txt
C:\Users\USER\AppData\Local\Temp\ps12-RUNID\top\sub\up
Last 3 paths:
C:\...\ps12-RUNID\top\sub\up\sub\up\sub ... \up\sub
C:\...\ps12-RUNID\top\sub\up\sub\up\sub ... \up\sub\file.txt
C:\...\ps12-RUNID\top\sub\up\sub\up\sub ... \up\sub\up
Maximum path characters:
505
\end{lstlisting}

\code{dir /s} follows the junction and lists the same three entries
(\code{sub}, \code{sub\textbackslash file.txt}, and \code{sub\textbackslash up})
again inside every \code{up}. Unlike GNU \code{ls}, it has no
check for a directory it has already listed. The output had 192 lines, which
is 3 entries per level for 64 levels. The longest path was 505 characters,
well past the classic 260-character \code{MAX\_PATH}, so path length was not
what stopped it.

The deepest paths in the transcript contain exactly 63 \code{\textbackslash up}
junctions. Windows allows at most 63 reparse points in one path. When
\code{dir} tries to open the 64th \code{up}, the open fails, so it cannot go
deeper and the listing ends. Windows does not detect the loop; it stops only
because of this limit.

\shot{windows/03b_lab1_directory_loop.png}{\linewidth}{Windows: junction loop
and \code{dir /s}.}

\paragraph{Answer.} Both systems allow links to folders only as symbolic links
(and junctions on Windows), never as hard links. A loop back to an upper-level
folder can be created on both. On Linux, \code{ls -R} is safe because it does
not follow links, and \code{ls -RL} and \code{find -L} detect the loop. On
Windows, \code{dir /s} follows the junction and repeats the same folders 63
times until it reaches the reparse-point limit.

\subsection{Link on drive D: to a source on drive C:}

A hard link is only a directory entry that names an inode or file record, and
those numbers are only meaningful inside one file system. A hard link
therefore cannot cross file systems or drives. A symbolic link stores a path,
so it can point anywhere.

\paragraph{Linux.} Linux has no drive letters, so the test used two separately
mounted file systems, \code{/tmp} and \code{/dev/shm}. Both are \code{tmpfs},
but they are different mounts with different device numbers.

\begin{lstlisting}
$ df -T /tmp /dev/shm | cat
Filesystem     Type  1K-blocks  Used Available Use% Mounted on
tmpfs          tmpfs   8120332 94276   8026056   2% /tmp
tmpfs          tmpfs   8120332 17944   8102388   1% /dev/shm
$ stat -c "%n device=%d" /tmp /dev/shm
/tmp device=55
/dev/shm device=28
$ ln /tmp/ps12-fsA.txt /dev/shm/ps12-hard.txt
ln: failed to create hard link '/dev/shm/ps12-hard.txt' => '/tmp/ps12-fsA.txt': Invalid cross-device link
$ ln -s /tmp/ps12-fsA.txt /dev/shm/ps12-soft.txt
$ cat /dev/shm/ps12-soft.txt
data on file system A
\end{lstlisting}

The hard link fails with \code{EXDEV} (``Invalid cross-device link''). The
symbolic link works.

\shot{linux/04_lab1_cross_filesystem.png}{0.85\linewidth}{Linux: links across
two file systems.}

\paragraph{Windows.}\mbox{}

\begin{lstlisting}
> mklink /H "D:\ps12-RUNID\hard.txt" "C:\Users\USER\AppData\Local\Temp\ps12-RUNID\hard.txt"
The system cannot move the file to a different disk drive.
> mklink "D:\ps12-RUNID\soft.txt" "C:\Users\USER\AppData\Local\Temp\ps12-RUNID\hard.txt"
symbolic link created for D:\ps12-RUNID\soft.txt <<===>> C:\...\hard.txt
> mklink /J "D:\ps12-RUNID\junction" "C:\Users\USER\AppData\Local\Temp\ps12-RUNID\top"
Junction created for D:\ps12-RUNID\junction <<===>> C:\...\top
> type "D:\ps12-RUNID\soft.txt"
original data
> type "D:\ps12-RUNID\junction\sub\file.txt"
file in sub
\end{lstlisting}

The hard link fails with ``The system cannot move the file to a different disk
drive.'' The file symbolic link and the directory junction on \code{D:} both
work, because both store a path to \code{C:}.

\shot{windows/04_lab1_cross_drive.png}{\linewidth}{Windows: links on D: to a
source on C:.}

\paragraph{Answer.} Yes for symbolic links and junctions; no for hard links.

\section{Lab 2: Named pipe (FIFO)}

\subsection{Create and list the pipe}

\begin{lstlisting}
$ mkfifo pipeFile
$ ls -la pipeFile
prw-r--r-- 1 gunteed gunteed 0 Nov 22 18:01 pipeFile
$ stat pipeFile
  File: pipeFile
  Size: 0         	Blocks: 0          IO Block: 4096   fifo
Device: 0,55	Inode: 3038        Links: 1
$ file pipeFile
pipeFile: fifo (named pipe)
\end{lstlisting}

The first character of the mode is \code{p}, which marks a named pipe. The
file has an inode and a name in the directory, but its size and block count
are 0. The pipe never stores data on disk.

A separate test shows that a writer blocks when no reader exists:

\begin{lstlisting}
$ timeout 2 sh -c 'echo hello > pipeFile'
[exit status 124]
\end{lstlisting}

Exit status 124 means \code{timeout} had to kill the command. The shell never
finished opening \code{pipeFile} for writing.

\shot{linux/05_lab2_fifo_setup.png}{0.75\linewidth}{Creating and listing the
named pipe.}

\subsection{Two terminals}

Two real terminals were driven with \code{tmux}. Terminal~1 (T1) is the writer
and Terminal~2 (T2) is the reader. Both screens were captured after each step.

\begin{lstlisting}
--- Step 1: T1 runs cat > pipeFile; no reader yet, so T1 waits in open() ---
[Terminal 1]
T1$ cat > pipeFile
[Terminal 2]
T2$
--- Step 2: T2 runs cat pipeFile; both ends are now open ---
[Terminal 1]
T1$ cat > pipeFile
[Terminal 2]
T2$ cat pipeFile
--- Step 3: T1 types hello linux ---
[Terminal 1]
T1$ cat > pipeFile
hello linux
[Terminal 2]
T2$ cat pipeFile
hello linux
--- Step 4: T1 types two more lines ---
[Terminal 1]
T1$ cat > pipeFile
hello linux
second line from terminal 1
third line
[Terminal 2]
T2$ cat pipeFile
hello linux
second line from terminal 1
third line
--- Step 5: T1 presses ^D; both cat commands exit ---
[Terminal 1]
...
third line
T1$
[Terminal 2]
...
third line
T2$
--- After the session ---
$ ls -la pipeFile
prw-r--r-- 1 gunteed gunteed 0 Nov 22 18:01 pipeFile
\end{lstlisting}

Each line typed in T1 appears in T2 right after Enter. After \code{\^{}D},
both commands return to the prompt. The pipe file is still there and its size
is still 0.

\shot{linux/06_lab2_fifo_terminals.png}{0.7\linewidth}{Two-terminal session
through the named pipe.}

\subsection{Concept and explanation}

A named pipe, or FIFO, is a pipe that has a name in the file system. An
ordinary pipe (\code{|}) exists only between a parent and its children,
because they inherit its file descriptors. A FIFO can be opened by name, so
unrelated processes, such as two shells in two terminals, can use it to
communicate. The file system entry is only a meeting point. The data passes
through a buffer in kernel memory (64~KiB by default on Linux), in first-in,
first-out order, and it is never written to disk. This is why the size is 0
before, during, and after the transfer.

The experiment behaves this way because of four rules:

\begin{enumerate}
    \item \textbf{Open blocks until both ends exist.} \code{open()} of a FIFO
    for writing waits until a process opens it for reading, and the reverse.
    In step 1, the shell in T1 is stuck opening \code{pipeFile} for the
    \code{>} redirection, so \code{cat} has not even started. In step 2, T2
    opens the read end and both opens return.
    \item \textbf{The terminal sends input one line at a time.} The terminal is
    in canonical mode, so \code{cat} in T1 receives a line from
    \code{read()} only after Enter. It writes that line into the pipe at once.
    \item \textbf{Read blocks until data arrives.} \code{cat} in T2 waits in
    \code{read()} on the empty pipe. When T1 writes a line, the kernel wakes
    T2, which prints the line. Each later line repeats the same cycle.
    \item \textbf{Closing the write end gives end-of-file.} \code{\^{}D} on
    an empty line makes \code{read()} on the terminal return 0, so \code{cat}
    in T1 exits and its write end of the pipe is closed. With no writers left,
    \code{read()} in T2 returns 0 once the buffer is empty. \code{cat} treats
    that as end-of-file and exits too.
\end{enumerate}

If the reader exited first, the next write in T1 would get \code{SIGPIPE} and
\code{cat} would be killed (``Broken pipe''). A FIFO is one-way; two-way
communication needs two FIFOs.

\section{Lab 3: Reading and writing files}

\subsection{Trace}

The file \code{foo} was created with \code{echo hello > foo}, and \code{cat foo}
was traced on a terminal. The lines before the \code{openat} of \code{foo}
load the C library and the locale and are omitted here.

\begin{lstlisting}
$ strace cat foo
...
fstat(1, {st_mode=S_IFCHR|0600, st_rdev=makedev(0x88, 0x7), ...}) = 0
openat(AT_FDCWD, "foo", O_RDONLY)       = 3
fstat(3, {st_mode=S_IFREG|0644, st_size=6, ...}) = 0
fadvise64(3, 0, 0, POSIX_FADV_SEQUENTIAL) = 0
mmap(NULL, 270336, PROT_READ|PROT_WRITE, MAP_PRIVATE|MAP_ANONYMOUS, -1, 0) = 0x7f12ac7be000
read(3, "hello\n", 262144)              = 6
write(1, "hello\n", 6hello
)                  = 6
read(3, "", 262144)                     = 0
close(3)                                = 0
munmap(0x7f12ac7be000, 270336)          = 0
close(1)                                = 0
close(2)                                = 0
exit_group(0)                           = ?
+++ exited with 0 +++
\end{lstlisting}

The four lines match the slide. The only difference is the buffer size: GNU
coreutils 9.12 asks for 262144 bytes (256~KiB), while the slide's system asked
for 32768 bytes (32~KiB). Newer coreutils use a larger I/O buffer to make fewer
system calls on large files. The \code{mmap} of 270336 bytes just before the
first \code{read} allocates that buffer.

\shot{linux/07_lab3_strace_key_calls.png}{0.85\linewidth}{\code{strace} of
\code{cat foo}, file-related calls only.}

\subsection{The four lines}

\begin{description}
    \item[\code{read(3, "hello\textbackslash n", 32768) = 6}] \code{cat} asks the
    kernel to read up to 32768 bytes from file descriptor 3, which is
    \code{foo}. The kernel copies the file content into the buffer and returns
    6, the number of bytes actually read. The file holds only 6 bytes, so the
    request is not filled.
    \item[\code{write(1, "hello\textbackslash n", 6) = 6}] \code{cat} writes
    those 6 bytes to file descriptor 1, standard output, which is the
    terminal. The return value 6 means all bytes were written. The text
    \code{6hello} appears because \code{strace} prints the start of the line
    to standard error before the call runs. The call then prints
    \code{hello} and a newline on the same terminal, and \code{strace} finishes
    the line with \code{) = 6} afterwards.
    \item[\code{read(3, "", 32768) = 0}] \code{cat} reads again. The file
    offset is already at the end, so the kernel returns 0. A return value of 0
    from \code{read()} means end-of-file, so \code{cat} stops its loop.
    \item[\code{close(3) = 0}] \code{cat} closes \code{foo}. The kernel frees
    descriptor 3 and its open file entry. The return value 0 means success.
\end{description}

\subsection{What the numbers represent}

\begin{center}
\begin{tabularx}{\linewidth}{@{}lX@{}}
\toprule
Number & Meaning \\
\midrule
\code{3} & The file descriptor for \code{foo}, returned by
\code{openat(AT\_FDCWD, "foo", O\_RDONLY) = 3}. Descriptors 0, 1, and 2 are
already open as standard input, output, and error. The kernel always returns
the lowest free number, which is 3. It is an index into the process's file
descriptor table, which points to the system-wide open file table entry that
holds the file offset and the inode. \\
\code{1} & The file descriptor for standard output. The shell set it up before
starting \code{cat}; here it is the terminal (\code{S\_IFCHR}, a character
device). \\
\code{6} & A byte count: \code{h e l l o} plus the newline \code{\textbackslash n}.
It is the return value of \code{read} (bytes read), the third argument of
\code{write} (bytes to write), and the return value of \code{write} (bytes
written). \\
\code{32768} & The size of \code{cat}'s buffer: the maximum number of bytes
that one \code{read()} may return. It is a request, not the file size. Our
system used 262144. \\
\bottomrule
\end{tabularx}
\end{center}

\subsection{Extra observation: output to a regular file}

When standard output is redirected to a regular file, this version of
\code{cat} does not use \code{read} and \code{write} at all:

\begin{lstlisting}
$ strace -e trace=openat,read,write,copy_file_range,close cat foo > out.txt
...
openat(AT_FDCWD, "foo", O_RDONLY)       = 3
copy_file_range(3, NULL, 1, NULL, 9223372035781033984, 0) = 6
copy_file_range(3, NULL, 1, NULL, 9223372035781033984, 0) = 0
close(3)                                = 0
\end{lstlisting}

\code{copy\_file\_range} asks the kernel to copy from descriptor 3 to
descriptor 1 directly, without moving the data through a buffer in user space.
The same numbers appear: 3 and 1 are the descriptors, 6 is the byte count, and
0 again means end-of-file. The traced calls depend on what the descriptors
refer to, not only on the command.

\section{Conclusion}

Hard links are extra names for one inode or file record, so they survive
deletion of any other name, but they cannot point to directories or cross file
systems. Symbolic links and junctions store a path, so they can point to
directories and other drives, but they break when the target name disappears
and they can create loops in the directory graph. GNU \code{ls} and
\code{find} detect those loops, while Windows \code{dir /s} follows them until
the 63-reparse-point limit.

A named pipe is a file system name for a kernel buffer. Blocking
\code{open}, blocking \code{read}, and end-of-file on writer close explain the
two-terminal behavior.

\code{strace} shows that \code{cat} works through file descriptors: it opens
the file as descriptor 3, reads with a fixed-size buffer, writes the byte
count it received to descriptor 1, and stops when \code{read} returns 0.

\end{document}
